\section{Fallback, Recovery, and Gap Adaptation}
\label{sec:s5g_gap}

The presence of transitional states is motivated by a very specific problem: an abrupt \emph{fallback} can be dynamically unsafe. If the platoon is traveling with a cooperative controller and with short gaps, an immediate switch to a more conservative controller does not automatically guarantee that the new safety conditions are already met at that instant. The new controller has safe gaps that differ from those of the previous one, and at the moment of the change it is therefore operating in an unsafe range \cite{ploeg2015graceful,segata2023multi-technology}.

For this reason, \emph{SafeSwitch} introduces a \emph{gap control} phase. In the event of degraded connectivity, the system does not move directly from a stable state to the next, more conservative stable state, but first enters an intermediate state in which the gap is progressively increased. Only when a safe distance for the new controller has been reached can the transition be considered complete. In the opposite direction, when link quality improves, recovery of the higher-performance state takes place by gradually reducing the gap before consolidating the new controller.

This choice is important because it logically separates two different problems: recognizing that connectivity is no longer sufficient to support a given level of cooperation, and bringing the platoon into a physical configuration that is consistent with the new safety level. The paper shows that a naive strategy based on a simple direct switch from CACC to ACC produces worse and unsafe behaviors in the presence of strong disturbances precisely because it lacks an intermediate adaptation phase \cite{segata2023multi-technology}.
